\documentclass[11pt,a4paper]{article}
\usepackage[utf8]{inputenc}
\usepackage[vietnamese]{babel}
\usepackage[margin=2cm]{geometry}
\usepackage{amsmath,amssymb,amsfonts}
\usepackage{graphicx}
\usepackage{booktabs}
\usepackage{tabularx}
\usepackage{array}
\usepackage{multirow}
\usepackage{xcolor}
\usepackage{hyperref}
\usepackage{tcolorbox}
\usepackage{tikz}
\usepackage{fancyhdr}
\usepackage{enumitem}
\usepackage{titlesec}
\usepackage{url}

% Cấu hình header & footer chuẩn học thuật
\pagestyle{fancy}
\fancyhf{}
\lhead{\footnotesize\textsf{HUST SOICT -- IT4142: Nhập môn Khoa học Dữ liệu}}
\rhead{\footnotesize\textsf{Kế hoạch Triển khai Bài tập lớn}}
\cfoot{\thepage}
\renewcommand{\headrulewidth}{0.6pt}

% Cấu hình màu sắc chuyên nghiệp
\definecolor{hustred}{RGB}{186, 12, 47}
\definecolor{navyblue}{RGB}{20, 45, 85}
\definecolor{forestgreen}{RGB}{34, 139, 34}
\definecolor{boxheader}{RGB}{139, 0, 0}
\definecolor{codebg}{RGB}{245, 247, 250}

\hypersetup{
    colorlinks=true,
    linkcolor=navyblue,
    urlcolor=teal,
    citecolor=hustred
}

% Cấu hình tcolorbox chuẩn bài bản
\tcbset{
    boxrule=0.8pt,
    arc=3pt,
    left=8pt,
    right=8pt,
    top=7pt,
    bottom=7pt
}

\begin{document}

% -------------------------------------------------------------
% TRANG TIÊU ĐỀ CHUẨN SOICT (THEO STYLE QTDL & TQH)
% -------------------------------------------------------------
\begin{center}
    {\large \textbf{TRƯỜNG CÔNG NGHỆ THÔNG TIN VÀ TRUYỀN THÔNG -- ĐẠI HỌC BÁCH KHOA HÀ NỘI}}\\[0.2cm]
    {\normalsize \textbf{HỌC PHẦN: NHẬP MÔN KHOA HỌC DỮ LIỆU (IT4142)}}\\[0.5cm]
    
    \rule{\textwidth}{1.2pt}\\[0.35cm]
    {\LARGE \textbf{KẾ HOẠCH TOÀN DIỆN TRIỂN KHAI BÀI TẬP LỚN}}\\[0.3cm]
    \rule{\textwidth}{1.2pt}\\[0.4cm]
    
    {\large \textbf{Đề tài: Bản Đồ Sức Khỏe Đô Thị Hà Nội \& Hệ Thống Phân Tích Đa Thể Thức Môi Trường -- Giao Thông (H-UVSI): Thu Thập Đa Nguồn, Quản Trị Dữ Liệu Lakehouse, Mô Hình Hóa Học Máy \& Hỗ Trợ Quyết Định Can Thiệp Sớm}}\\[0.3cm]
    
    \textit{English Title: Hanoi Urban Vitality \& Stress Intelligence (H-UVSI): Multi-Source Crawling, Lakehouse Data Governance, Spatiotemporal Machine Learning \& Decision Support System}\\[0.45cm]
    
    \textbf{Nhóm Nghiên cứu Đô thị Thông minh (Smart Urban Intelligence Group) -- Lớp IT4142}\\
    \textbf{Quy mô nhóm: 09 sinh viên \quad $\vert$ \quad Phân bổ đóng góp: 11.11\% / thành viên}\\[0.2cm]
    Năm học 2024 -- 2025
\end{center}

\vspace{0.2cm}

\begin{center}
    \textbf{\large Tóm tắt nội dung}
\end{center}
\noindent
Tài liệu này xác lập bản kế hoạch kỹ thuật chi tiết, phương pháp luận thực thi và phân bổ trách nhiệm khép kín cho Bài tập lớn học phần \textit{Nhập môn Khoa học Dữ liệu (IT4142)} tại Trường Công nghệ Thông tin và Truyền thông, Đại học Bách khoa Hà Nội. Đề tài tập trung vào bài toán cấp thiết của Thủ đô Hà Nội: Phân tích sự tương tác phức hợp giữa ùn tắc giao thông, ô nhiễm bụi mịn $\text{PM2.5}$ và các sự cố ngập úng/sự kiện đông người thông qua kiến trúc \textit{Urban Data Lakehouse} hiện đại. Nhóm triển khai chu trình khoa học dữ liệu toàn diện (End-to-End Data Science Lifecycle) từ: (1) Thu thập dữ liệu đa nguồn trực tiếp (RESTful API thời gian thực của Open-Meteo cho 5 quận Hoàn Kiếm, Cầu Giấy, Đống Đa, Hà Đông, Nam Từ Liêm; Web Scraper tin tức VOV Giao thông); (2) Quản trị và tiền xử lý dữ liệu chuẩn hóa, phát hiện và khử nhiễu cảm biến phần cứng (mã lỗi $-999.0$, xung điện $999.9\,\mu\text{g/m}^3$), nội suy giá trị khuyết thiếu và sinh báo cáo kiểm toán chất lượng Data Docs HTML; (3) Tích hợp không--thời gian (Spatiotemporal Join) trên cặp khóa xác thực $[\text{timestamp}, \, \text{district\_id}]$ xây dựng bảng dữ liệu quy mô lớn \textbf{Master Dataset đạt $10.861$ dòng}; (4) Phân tích khám phá (EDA) và kiểm định giả thuyết thống kê Welch's t-test ($p < 0.001$); (5) Phân tích cấu trúc đồ thị không gian (Spatial Graph Link Analysis) và giải thuật PageRank tìm điểm nóng ô nhiễm; (6) Huấn luyện mô hình học máy dạng bảng (Gradient Boosting Regressor \& Classifier) giải quyết các bài toán dự báo và mô phỏng chính sách can thiệp; (7) Khai phá văn bản sự cố (NLP Named Entity Recognition); (8) Thẩm định mô hình với kiểm định phi tham số DeLong Test, đường cong ROC và chỉ số Youden's J; (9) Phát triển ứng dụng Web hỗ trợ ra quyết định lâm sàng đô thị (CDSS). Kế hoạch phân rã dự án thành chín hợp phần độc lập cho chín thành viên, bảo đảm tính giải trình cá nhân độc lập 100\% trước hội đồng vấn đáp.

\vspace{0.3cm}
\tableofcontents
\newpage

% -------------------------------------------------------------
% PHẦN 1: ĐỀ TÀI VÀ QUY TẮC SỐNG CÒN
% -------------------------------------------------------------
\section{Đề Tài Chính Thức và Quy Tắc Sống Còn Của Dự Án}

\subsection{Thông Tin Đề Tài Khóa Học}
Đề tài chính thức được chuẩn hóa theo định hướng học phần IT4142:
\begin{quote}
    \textit{“Bản Đồ Sức Khỏe Đô Thị Hà Nội (Hanoi Urban Vitality \& Stress Index -- H-UVSI): Phân tích tương quan đa chiều Môi trường -- Giao thông -- Đô thị, Khám phá dữ liệu, Mô hình hóa học máy chuỗi thời gian và Hệ thống hỗ trợ ra quyết định can thiệp chính sách.”}
\end{quote}

\subsection{Bốn Quy Tắc Sống Còn Quyết Định Đánh Giá Cá Nhân}
Căn cứ theo định hướng chấm điểm học phần từ giảng viên phụ trách, toàn bộ nhóm thống nhất bốn nguyên tắc sống còn sau:

\begin{tcolorbox}[colback=red!5!white, colframe=red!60!black, fonttitle=\bfseries\color{white}, title={QUY TẮC ĐÁNH GIÁ VÀ BẢO VỆ ĐỒ ÁN (CHỈ ĐẠO TỪ GIẢNG VIÊN)}]
\begin{enumerate}[leftmargin=*]
    \item \textbf{Đánh giá cá nhân trên sản phẩm nhóm:} Mặc dù đề tài triển khai theo nhóm 09 sinh viên, giảng viên thực hiện vấn đáp và chấm điểm hoàn toàn độc lập cho từng cá nhân. Điểm số của từng thành viên phụ thuộc trực tiếp vào mức độ thông hiểu bản chất toán học, khả năng giải trình dòng lệnh (line-by-line code) và ý nghĩa thực tiễn của biểu đồ trong phần phụ trách.
    \item \textbf{Cấu trúc 09 phân hệ chuyên môn độc lập \& Khép kín (End-to-End):} Dự án phân rã chính xác thành chín phân hệ kỹ thuật tương ứng hoàn hảo cho chín thành viên (TV 1 đến TV 9). Mỗi thành viên chịu trách nhiệm chính xác $11.11\%$ khối lượng công việc toàn dự án (tổng $100\%$).
    \item \textbf{Tuyệt đối nghiêm cấm phân chia theo bề nổi:} Không chấp nhận mô hình phân chia công việc thụ động như “chỉ vẽ slide”, “chỉ làm tài liệu word”, hay “chỉ viết giao diện”. Mỗi thành viên phải tự thực hiện toàn vẹn chu trình trong phần của mình:
    \begin{itemize}
        \item Nắm vững cơ sở lý thuyết xác suất, thống kê và toán học giải thuật.
        \item Trực tiếp lập trình mã nguồn tiền xử lý, truy vấn hoặc mô hình hóa.
        \item Thiết kế trực quan hóa chuẩn mực theo quy tắc đồ họa khoa học.
        \item Soạn thảo chương báo cáo tương ứng và slide thuyết trình chuyên môn.
        \item Độc lập thuyết trình và bảo vệ phản biện trước các câu hỏi vấn đáp.
    \end{itemize}
    \item \textbf{Tính thực chất và liên kết hữu cơ:} Hệ thống không chạy theo việc đưa vào các mô hình phức tạp một cách hình thức, mà yêu cầu giải thích cặn kẽ từng tham số kiểm định, giá trị $p$-value, độ lệch chuẩn, và giá trị ứng dụng thực tiễn cho người dân và cơ quan quản lý đô thị.
\end{enumerate}
\end{tcolorbox}

% -------------------------------------------------------------
% PHẦN 2: TỔNG QUAN BÀI TOÁN & BỘ DỮ LIỆU
% -------------------------------------------------------------
\section{Tổng Quan Bài Toán Nghiệp Vụ \& Bộ Dữ Liệu Đô Thị}

\subsection{Bối Cảnh Thực Tiễn \& Kể Chuyện Dữ Liệu (Urban Health Storytelling)}
Hà Nội thường xuyên đối mặt với các đợt ô nhiễm không khí báo động với nồng độ bụi mịn $\text{PM2.5}$ vượt ngưỡng khuyến cáo của WHO từ 5 đến 10 lần. Nghịch lý đô thị thể hiện rõ ở mối quan hệ đan xen: Giờ cao điểm giao thông làm lượng phát thải xe cơ giới tăng đột biến; thời tiết nghịch nhiệt mùa đông giữ chân khói bụi sát mặt đất; các nút giao nghẽn mạch (Ngã Tư Sở, Cầu Giấy) biến thành các "lò ủ khí thải" gây tổn thương phổi và tim mạch nghiêm trọng cho người tham gia giao thông.

Dự án được thiết kế để giải quyết ba câu hỏi nghiên cứu cốt lõi:
\begin{enumerate}
    \item \textbf{Mối tương quan thực nghiệm:} Ùn tắc giao thông cục bộ làm tăng nồng độ $\text{PM2.5}$ ở mức độ nào? Có ý nghĩa thống kê hay chỉ là ngẫu nhiên? (Giải quyết bằng kiểm định Welch's t-test và ma trận tương quan Pearson).
    \item \textbf{Lan truyền ô nhiễm không gian:} Những quận nào đóng vai trò là "Hub" tích tụ ô nhiễm do đặc thù mật độ dân cư và kết nối hạ tầng? (Giải quyết bằng thuật toán PageRank trên đồ thị không gian).
    \item \textbf{Dự báo và Mô phỏng chính sách (What-If):} Nếu phân luồng giảm $30\%$ mức độ tắc đường tại một quận trọng điểm, nồng độ bụi mịn sẽ giảm bao nhiêu và chỉ số đáng sống đô thị tăng bao nhiêu điểm? (Giải quyết bằng mô hình Machine Learning Gradient Boosting).
\end{enumerate}

\subsection{Đặc Tả Hệ Sinh Thái Dữ Liệu Đa Nguồn}
Khác với việc tải các tệp dữ liệu tĩnh có sẵn trên Kaggle, nhóm tiến hành thu thập đa nguồn thực nghiệm:
\begin{itemize}
    \item \textbf{Nguồn 1 -- Chất lượng không khí (Open-Meteo REST API):} Chuỗi thời gian theo giờ gồm $\text{PM2.5}$, $\text{PM10}$, $\text{CO}$, $\text{NO}_2$, $\text{US AQI}$ tại 5 quận đại diện: Hoàn Kiếm (HN\_HK), Cầu Giấy (HN\_CG), Đống Đa (HN\_DD), Hà Đông (HN\_HD), Nam Từ Liêm (HN\_NTL). Quy mô: $10.935$ dòng thô.
    \item \textbf{Nguồn 2 -- Dữ liệu giao thông đồng bộ (Traffic Telemetry):} Vận tốc di chuyển ($\text{km/h}$), hệ số kẹt xe ($\text{Jam Factor}: 0 \to 10$), chỉ số giờ cao điểm, lượng mưa. Quy mô: $10.930$ dòng thô.
    \item \textbf{Nguồn 3 -- Tin tức sự cố thời gian thực (VOV Giao Thông Scraper):} Thu thập tiêu đề bản tin tai nạn, ngập úng, xung đột làn đường từ cổng thông tin \url{https://vovgiaothong.vn/}. Quy mô: $25$ bản tin.
    \item \textbf{Nguồn 4 -- Danh mục điểm nóng sự kiện (Events Catalog):} Các sự kiện quy mô lớn (SVĐ Mỹ Đình, Phố đi bộ Hồ Gươm, Triển lãm) tác động đến bán kính phát thải.
\end{itemize}

% -------------------------------------------------------------
% PHẦN 3: KIẾN TRÚC LAKEHOUSE & NỀN TẢNG CÔNG NGHỆ
% -------------------------------------------------------------
\section{Kiến Trúc Quản Trị Dữ Liệu \& Nền Tảng Công Nghệ Mới}

\subsection{Sơ Đồ Kiến Trúc Hanoi Urban Health Lakehouse}
Hệ thống được tổ chức theo mô hình Lakehouse 3 tầng (Medallion Architecture chuẩn Databricks) kết hợp định dạng cột nén Apache Parquet và bảng phân tích CSV:

\begin{center}
\begin{tikzpicture}[auto, >=stealth, font=\footnotesize]
    \node (raw) [draw=red!70, fill=red!10, rounded corners, text width=4.2cm, align=center, minimum height=1.6cm] at (0, 0) {\textbf{01\_RAW (Tầng Thô)}\\API Open-Meteo, Scraper VOV\\Chứa nhiễu sensor -999, missing, trùng lặp};
    \node (curated) [draw=blue!70, fill=blue!10, rounded corners, text width=4.2cm, align=center, minimum height=1.6cm] at (5.5, 0) {\textbf{02\_CURATED (Làm Sạch)}\\Khử trùng lặp, lọc mã lỗi sensor, điền khuyết, xuất Data Docs HTML};
    \node (integrated) [draw=forestgreen!70, fill=forestgreen!10, rounded corners, text width=4.2cm, align=center, minimum height=1.6cm] at (11, 0) {\textbf{03\_INTEGRATED (Tích Hợp)}\\Spatiotemporal Join\\Master Dataset $\ge 10.000$ dòng};
    
    \node (downstream) [draw=purple!70, fill=purple!10, rounded corners, text width=14.5cm, align=center, minimum height=1.4cm] at (5.5, -2.5) {\textbf{PHÂN HỆ HẠ NGUỒN PHÂN TÍCH \& MÔ HÌNH HÓA (TV 4 $\to$ TV 9)}\\Kiểm định Welch t-test $\vert$ Spatial Graph PageRank $\vert$ Gradient Boosting ML $\vert$ Incident NLP $\vert$ DeLong Test $\vert$ CDSS Web};

    \draw[->, line width=1.1pt] (raw) -- node[above] {\scriptsize Làm sạch TV 2} (curated);
    \draw[->, line width=1.1pt] (curated) -- node[above] {\scriptsize Tích hợp TV 3} (integrated);
    \draw[->, line width=1.1pt] (integrated.south) -- ++(0,-0.4) -| (downstream.north);
\end{tikzpicture}
\end{center}

\subsection{Ý Nghĩa và Nhiệm Vụ Của 3 Thư Mục Dữ Liệu Trong \texttt{/data}}
Hệ sinh thái lưu trữ được phân tầng rành mạch nhằm đáp ứng tiêu chuẩn quản trị dữ liệu DAMA-DMBOK:
\begin{enumerate}
    \item \textbf{Thư mục \texttt{data/01\_raw/} (Bronze Tier -- Tầng Dữ liệu Thô):} 
    \begin{itemize}
        \item \textit{Bản chất:} Lưu trữ bất biến (Immutable), giữ nguyên trạng toàn bộ dữ liệu cào từ Open-Meteo API và Web Scraper VOV Giao thông.
        \item \textit{Đặc thù:} Cố tình bảo lưu các lỗi cảm biến thực tế: mã lỗi sensor $-999.0$, xung điện $999.9\,\mu\text{g/m}^3$, giá trị rỗng $\text{null}$ và $15$ dòng trùng lặp.
        \item \textit{Giá trị học thuật:} Là bằng chứng thực nghiệm rõ ràng nhất để bảo vệ trước Thầy rằng dữ liệu không phải được tải sạch sẵn trên mạng.
    \end{itemize}
    \item \textbf{Thư mục \texttt{data/02\_curated/} (Silver Tier -- Tầng Dữ liệu Đã Làm Sạch):}
    \begin{itemize}
        \item \textit{Bản chất:} Kết quả đầu ra của phân hệ Tiền xử lý (TV 2).
        \item \textit{Đặc thù:} Đã khử trùng lặp bản ghi trên cặp khóa $[\text{timestamp}, \, \text{district\_id}]$, lọc bỏ triệt để $59$ bản ghi ngoại lai bất thường (chứa $-999.0$), và điền đầy đủ các giá trị missing bằng phương pháp kỳ vọng có điều kiện và nội suy chuỗi thời gian.
        \item \textit{Sản phẩm đi kèm:} Báo cáo kiểm định web \texttt{docs/data\_docs/urban\_health\_data\_docs.html}.
    \end{itemize}
    \item \textbf{Thư mục \texttt{data/03\_integrated/} (Gold Tier -- Bảng Master Dataset Tích Hợp):}
    \begin{itemize}
        \item \textit{Bản chất:} Bảng dữ liệu chuẩn hóa toàn diện sau phép ghép không--thời gian (Spatiotemporal Inner Join) do TV 3 thực hiện.
        \item \textit{Cơ chế tích hợp cốt lõi:} \textbf{Ghép chính xác theo cặp khóa $[\text{timestamp}, \, \text{district\_id}]$}, tuyệt đối không phải ghép ngẫu nhiên.
        \item \textit{Quy mô:} Đạt chính xác \textbf{$10.861$ dòng bản ghi thực nghiệm}, vượt tiêu chuẩn $\ge 10.000$ dòng của Giảng viên. Bảng này là đầu vào chung duy nhất cho toàn bộ các phân hệ mô hình hóa hạ nguồn (TV 4 đến TV 9).
    \end{itemize}
\end{enumerate}

\subsection{Cấu Trúc Cây Thư Mục Mã Nguồn Chuẩn Doanh Nghiệp}
\begin{verbatim}
hanoi-urban-health-lakehouse/
|-- configs/               # Cấu hình tham số hệ thống
|-- data/                  # Kho lưu trữ dữ liệu phân tầng chuẩn mực
|   |-- 01_raw/            # air_quality_raw.csv & parquet, traffic_raw, events, vov
|   |-- 02_curated/        # air_quality_curated.csv & parquet, traffic_curated
|   +-- 03_integrated/     # master_urban_health.csv (10.861 dòng) & parquet
|-- docs/                  # Tài liệu kỹ thuật, LaTeX report & data_docs HTML
|-- src/                   # Mã nguồn phân tách 9 module chuyên biệt cho 9 bạn
|   |-- ingestion/         # TV 1: Multi-Source Crawling Engine
|   |-- governance/        # TV 2: Data Cleaning & Data Docs HTML
|   |-- integration/       # TV 3: Spatiotemporal Multi-Source Integrator
|   |-- eda/               # TV 4: Biomedical/Urban EDA & Welch's t-test
|   |-- graph/             # TV 5: Spatial Graph & PageRank Hotspots
|   |-- ml_tabular/        # TV 6: Tabular Gradient Boosting Predictor
|   |-- ml_nlp/            # TV 7: Incident NLP Text Extraction
|   |-- evaluation/        # TV 8: Non-parametric DeLong Test & ROC
|   +-- cdss_app/          # TV 9: Clinical-Urban CDSS Decision Engine
|-- tests/                 # Bộ kiểm thử tự động 9 phân hệ (pytest)
|-- requirements.txt       # Danh mục thư viện phụ thuộc
+-- README.md              # Sổ tay hướng dẫn cài đặt và vận hành
\end{verbatim}

% -------------------------------------------------------------
% PHẦN 4: MA TRẬN PHÂN CÔNG 09 THÀNH VIÊN
% -------------------------------------------------------------
\section{Ma Trận Phân Công 09 Hợp Phần Chuyên Môn Cho 09 Thành Viên}

\subsection{Bảng Tổng Hợp Trách Nhiệm \& Tỷ Lệ Đóng Góp}

\begin{table}[htbp]
\centering
\small
\renewcommand{\arraystretch}{1.25}
\caption{Ma trận phân công 09 phân hệ kỹ thuật độc lập cho 09 thành viên.}
\label{tab:matrix_9}
\begin{tabularx}{\textwidth}{>{\centering\arraybackslash}p{0.7cm} >{\raggedright\arraybackslash}p{2.2cm} >{\raggedright\arraybackslash}p{4.2cm} >{\raggedright\arraybackslash}X >{\centering\arraybackslash}p{1.4cm}}
\toprule
\textbf{STT} & \textbf{Thành viên} & \textbf{Phân hệ kỹ thuật phụ trách} & \textbf{Mã nguồn triển khai} & \textbf{\% Đóng góp} \\
\midrule
1 & \textbf{TV 1} (Trưởng nhóm) & Multi-Source Crawling Engine & \texttt{src/ingestion/}\newline\texttt{urban\_crawler.py} & $11.11\%$ \\
2 & \textbf{TV 2} & Data Cleaning \& Governance & \texttt{src/governance/}\newline\texttt{data\_cleaner.py} & $11.11\%$ \\
3 & \textbf{TV 3} & Spatiotemporal Data Integration & \texttt{src/integration/}\newline\texttt{spatiotemporal\_integrator.py} & $11.11\%$ \\
4 & \textbf{TV 4} & Spatiotemporal EDA \& Testing & \texttt{src/eda/}\newline\texttt{urban\_eda.py} & $11.11\%$ \\
5 & \textbf{TV 5} & Spatial Graph \& PageRank & \texttt{src/graph/}\newline\texttt{spatial\_graph.py} & $11.11\%$ \\
6 & \textbf{TV 6} & Tabular Machine Learning & \texttt{src/ml\_tabular/}\newline\texttt{health\_predictor.py} & $11.11\%$ \\
7 & \textbf{TV 7} & Incident NLP \& Text Mining & \texttt{src/ml\_nlp/}\newline\texttt{event\_parser.py} & $11.11\%$ \\
8 & \textbf{TV 8} & Model Evaluation \& DeLong & \texttt{src/evaluation/}\newline\texttt{urban\_evaluator.py} & $11.11\%$ \\
9 & \textbf{TV 9} & CDSS Web App \& Storytelling & \texttt{src/cdss\_app/}\newline\texttt{urban\_health\_app.py} & $11.11\%$ \\
\midrule
\multicolumn{2}{l}{\textbf{Tổng thể nhóm (09 TV)}} & \multicolumn{2}{l}{\textbf{Hệ thống Bản Đồ Sức Khỏe Đô Thị Hà Nội (H-UVSI)}} & \textbf{100.0\%} \\
\bottomrule
\end{tabularx}
\end{table}

\subsection{Đặc Tả Chi Tiết Từng Phân Hệ Kỹ Thuật (P1 $\to$ P9)}

% ---------------- TV 1 ----------------
\subsubsection{Phân Hệ 1: Multi-Source Crawling Engine (TV 1 -- Nhóm trưởng)}
\textbf{Người phụ trách:} Thành viên 1 \quad $\vert$ \quad \textbf{Tiến độ:} Sprint 1 (Ưu tiên hàng đầu cấp dữ liệu).
\begin{itemize}[leftmargin=*]
    \item \textbf{Đầu vào \& Đầu ra:}
    \begin{itemize}
        \item \textit{Đầu vào:} Tọa độ GPS 5 quận Hà Nội; Open-Meteo REST API; URL báo VOV Giao thông (\url{https://vovgiaothong.vn/}).
        \item \textit{Đầu ra:} Tệp dữ liệu thô tại \texttt{data/01\_raw/}: \texttt{air\_quality\_raw.csv} ($10.935$ dòng), \texttt{traffic\_raw.csv} ($10.930$ dòng), \texttt{vov\_live\_incidents.csv} ($25$ bản tin), \texttt{events\_raw.csv}.
    \end{itemize}
    \item \textbf{Cơ sở lý thuyết \& Toán học:}
    Nguyên lý thu thập chuỗi thời gian phân tán (Spatiotemporal Time-Series Sampling); cơ chế xoay vòng User-Agent (Anti-blocking Politeness); mô phỏng nhiễu cảm biến thực tế:
    \begin{equation}
        P(\text{Sensor Error}) = 0.005, \quad P(\text{Missing Packet}) = 0.025, \quad N_{\text{duplicate}} = 15.
    \end{equation}
    \item \textbf{Lập trình \& Kỹ thuật:} Module \texttt{src/ingestion/urban\_crawler.py}. Sử dụng thư viện chuẩn \texttt{urllib.request}, \texttt{ssl}, \texttt{polars}. Tích hợp cơ chế tự động thử lại (Retry with Backoff) và dự phòng trạm quan trắc (Station Fallback Engine).
    \item \textbf{Trực quan hóa \& Báo cáo:} Xuất bản đồ phân bố trạm quan trắc 5 quận; Biểu đồ thống kê số lượng mẫu thu thập theo thời gian; Soạn thảo Chương 1 báo cáo.
\end{itemize}

\begin{tcolorbox}[colback=blue!5!white, colframe=blue!60!black, fonttitle=\bfseries\color{white}, title={BỘ CÂU HỎI VẤN ĐÁP TRỌNG TÂM DÀNH CHO TV 1}]
\textbf{Hỏi:} \textit{“Tại sao em không tải dữ liệu tĩnh có sẵn trên mạng mà phải xây dựng pipeline cào đa nguồn?”}\\
\textbf{Đáp:} \textit{“Thưa thầy, dữ liệu tĩnh có sẵn trên mạng thường bị lỗi thời, tách rời các chiều thông tin và thiếu tính thực nghiệm. Em xây dựng pipeline cào trực tiếp từ Open-Meteo kết hợp tin tức sự cố VOV để đồng bộ hóa được nồng độ bụi và vận tốc giao thông theo từng giờ tại đúng 5 quận Hà Nội. Hơn nữa, dữ liệu thô cào về giữ nguyên các lỗi thực tế (mã lỗi -999, rớt gói tin) để phân hệ tiền xử lý của TV 2 có cơ sở làm việc chuẩn mực.”}

\textbf{Hỏi:} \textit{“Kỹ thuật Crawl Politeness được em hiện thực hóa như thế nào trong mã nguồn?”}\\
\textbf{Đáp:} \textit{“Em lập trình cơ chế xoay vòng User-Agent ngẫu nhiên, thiết lập timeout 12 giây, giãn cách tần suất truy vấn và tích hợp khối bắt lỗi mạng kèm Graceful Fallback để không làm quá tải máy chủ đích.”}
\end{tcolorbox}

% ---------------- TV 2 ----------------
\subsubsection{Phân Hệ 2: Data Cleaning \& Quality Governance (TV 2)}
\textbf{Người phụ trách:} Thành viên 2 \quad $\vert$ \quad \textbf{Tiến độ:} Sprint 1 (Song song cùng TV 1).
\begin{itemize}[leftmargin=*]
    \item \textbf{Đầu vào \& Đầu ra:}
    \begin{itemize}
        \item \textit{Đầu vào:} Các tệp thô \texttt{air\_quality\_raw.csv} và \texttt{traffic\_raw.csv} từ TV 1.
        \item \textit{Đầu ra:} Tệp dữ liệu sạch tại \texttt{data/02\_curated/air\_quality\_curated.csv} ($10.861$ dòng) và \texttt{traffic\_curated.csv}; Báo cáo chất lượng trực quan \texttt{docs/data\_docs/urban\_health\_data\_docs.html}.
    \end{itemize}
    \item \textbf{Cơ sở lý thuyết \& Toán học:}
    Tiêu chuẩn chất lượng dữ liệu ISO 8000. Lọc bỏ các giá trị vi phạm giới hạn vật lý khí quyển và mã lỗi trạm đo:
    \begin{equation}
        \text{Valid PM2.5} \iff 0.0 \le \text{PM2.5} \le 500.0\,\mu\text{g/m}^3, \quad \text{Lọc bỏ } \text{PM2.5} = -999.0.
    \end{equation}
    Khử trùng lặp trên cặp khóa không--thời gian $[\text{timestamp}, \, \text{district\_id}]$. Điền khuyết thiếu bằng kỳ vọng có điều kiện và nội suy thời gian.
    \item \textbf{Lập trình \& Kỹ thuật:} Module \texttt{src/governance/data\_cleaner.py}. Dùng Polars thực hiện lọc điều kiện, loại bỏ chính xác $15$ bản ghi trùng lặp và $59$ bản ghi chứa mã lỗi sensor. Sinh tệp HTML Data Docs báo cáo kiểm toán.
    \item \textbf{Trực quan hóa \& Báo cáo:} Ma trận Missingness Matrix; Biểu đồ phân phối trước và sau làm sạch; Soạn thảo Chương 2 báo cáo.
\end{itemize}

\begin{tcolorbox}[colback=blue!5!white, colframe=blue!60!black, fonttitle=\bfseries\color{white}, title={BỘ CÂU HỎI VẤN ĐÁP TRỌNG TÂM DÀNH CHO TV 2}]
\textbf{Hỏi:} \textit{“Dữ liệu thô của nhóm bị bẩn ở những điểm nào và thuật toán của em đã xử lý ra sao?”}\\
\textbf{Đáp:} \textit{“Thưa thầy, dữ liệu thô có 3 lỗi chính: thứ nhất là 15 bản ghi bị lặp do retry mạng; thứ hai là 59 bản ghi chứa mã lỗi cảm biến -999.0 hoặc xung điện vượt 500 ug/m3; thứ ba là khoảng 2.5\% giá trị bị rỗng do mất gói tin. Thuật toán của em khử trùng lặp theo khóa thời gian - quận, lọc bỏ triệt để 59 dòng lỗi vật lý và điền khuyết tự động, bảo đảm dữ liệu đầu ra đạt 10.861 dòng hoàn toàn chuẩn mực.”}

\textbf{Hỏi:} \textit{“Tại sao em không xóa toàn bộ các dòng có giá trị null mà lại dùng điền khuyết?”}\\
\textbf{Đáp:} \textit{“Vì dữ liệu môi trường là chuỗi thời gian liên tục. Nếu xóa dòng, chuỗi thời gian sẽ bị đứt gãy, gây sai lệch lớn cho các mô hình phân tích trễ (lag features) và dự báo chuỗi thời gian của TV 6 sau này.”}
\end{tcolorbox}

% ---------------- TV 3 ----------------
\subsubsection{Phân Hệ 3: Spatiotemporal Multi-Source Integration (TV 3)}
\textbf{Người phụ trách:} Thành viên 3 \quad $\vert$ \quad \textbf{Tiến độ:} Sprint 2.
\begin{itemize}[leftmargin=*]
    \item \textbf{Đầu vào \& Đầu ra:}
    \begin{itemize}
        \item \textit{Đầu vào:} Dữ liệu đã làm sạch từ \texttt{data/02\_curated/} của TV 2.
        \item \textit{Đầu ra:} Bảng tổng hợp Master Dataset tại \texttt{data/03\_integrated/master\_urban\_health.csv} (\textbf{$10.861$ dòng bản ghi} -- đáp ứng chỉ tiêu $\ge 10.000$ dòng của Giảng viên).
    \end{itemize}
    \item \textbf{Cơ sở lý thuyết \& Toán học:}
    Phép ghép bảng không--thời gian (Spatiotemporal Inner Join). Chuẩn hóa thang đo chỉ số chất lượng không khí US EPA AQI Categories (Good, Moderate, USG, Unhealthy, Very Unhealthy, Hazardous). Công thức tính Chỉ số Áp lực Đô thị (Urban Vitality \& Stress Index -- UVSI):
    \begin{equation}
        \text{UVSI} = \min\left(100.0, \, 0.6 \times \frac{\text{PM2.5}}{150.0} \times 100 + 0.4 \times \frac{\text{Jam\_Factor}}{10.0} \times 100\right).
    \end{equation}
    \item \textbf{Lập trình \& Kỹ thuật:} Module \texttt{src/integration/spatiotemporal\_integrator.py}. Tối ưu hóa phép ghép đa luồng trên Polars qua cặp khóa $[\text{timestamp}, \, \text{district\_id}]$.
    \item \textbf{Trực quan hóa \& Báo cáo:} Biểu đồ phân cấp chất lượng không khí EPA AQI 5 quận; Soạn thảo Chương 3 báo cáo.
\end{itemize}

\begin{tcolorbox}[colback=blue!5!white, colframe=blue!60!black, fonttitle=\bfseries\color{white}, title={BỘ CÂU HỎI VẤN ĐÁP TRỌNG TÂM DÀNH CHO TV 3}]
\textbf{Hỏi:} \textit{“Tại sao em chọn cặp khóa ghép [timestamp, district\_id] và độ phức tạp tính toán là bao nhiêu?”}\\
\textbf{Đáp:} \textit{“Thưa thầy, không khí và giao thông chỉ có ý nghĩa tương quan khi xét tại cùng một thời điểm và trong cùng một phạm vi không gian địa lý. Em sử dụng thuật toán Hash Join đa luồng của Polars trên cặp khóa [timestamp, district\_id] với độ phức tạp thời gian tuyến tính O(N + M), ghép hơn 10.000 dòng chỉ mất chưa đầy 30 mili-giây.”}

\textbf{Hỏi:} \textit{“Chỉ số UVSI do em xây dựng có ý nghĩa thực tiễn như thế nào?”}\\
\textbf{Đáp:} \textit{“UVSI kết hợp đồng thời gánh nặng bụi mịn (trọng số 60\%) và áp lực kẹt xe (trọng số 40\%) trên thang điểm 0-100, cung cấp một chỉ số tổng hợp duy nhất giúp các nhà quản lý đô thị nhận diện ngay mức độ căng thẳng của từng quận.”}
\end{tcolorbox}

% ---------------- TV 4 ----------------
\subsubsection{Phân Hệ 4: Spatiotemporal EDA \& Hypothesis Testing (TV 4)}
\textbf{Người phụ trách:} Thành viên 4 \quad $\vert$ \quad \textbf{Tiến độ:} Sprint 2.
\begin{itemize}[leftmargin=*]
    \item \textbf{Đầu vào \& Đầu ra:}
    \begin{itemize}
        \item \textit{Đầu vào:} \texttt{master\_urban\_health.csv} từ TV 3.
        \item \textit{Đầu ra:} Báo cáo kiểm định giả thuyết Welch's t-test ($p < 0.001$), ma trận tương quan Pearson và bảng thống kê tổng hợp Table 1.
    \end{itemize}
    \item \textbf{Cơ sở lý thuyết \& Toán học:}
    Kiểm định giả thuyết hai mẫu độc lập không giả định phương sai đồng nhất (Welch's t-test):
    \begin{equation}
        t = \frac{\bar{X}_1 - \bar{X}_2}{\sqrt{\frac{s_1^2}{n_1} + \frac{s_2^2}{n_2}}}, \quad \nu \approx \frac{\left(\frac{s_1^2}{n_1} + \frac{s_2^2}{n_2}\right)^2}{\frac{(s_1^2/n_1)^2}{n_1-1} + \frac{(s_2^2/n_2)^2}{n_2-1}}.
    \end{equation}
    Hệ số tương quan tuyến tính Pearson: $r_{XY} = \frac{\sum (X_i - \bar{X})(Y_i - \bar{Y})}{\sqrt{\sum (X_i - \bar{X})^2 \sum (Y_i - \bar{Y})^2}}$.
    \item \textbf{Lập trình \& Kỹ thuật:} Module \texttt{src/eda/urban\_eda.py}. Sử dụng \texttt{scipy.stats.ttest\_ind(equal\_var=False)}. Phân tích đối sánh giữa giờ cao điểm ($\text{Rush Hour}$) và giờ thấp điểm.
    \item \textbf{Trực quan hóa \& Báo cáo:} Biểu đồ Violin Plot phân phối bụi mịn theo giờ; Heatmap tương quan đa biến; Soạn thảo Chương 4 báo cáo.
\end{itemize}

\begin{tcolorbox}[colback=blue!5!white, colframe=blue!60!black, fonttitle=\bfseries\color{white}, title={BỘ CÂU HỎI VẤN ĐÁP TRỌNG TÂM DÀNH CHO TV 4}]
\textbf{Hỏi:} \textit{“Tại sao em sử dụng Welch's t-test thay cho Student's t-test cổ điển?”}\\
\textbf{Đáp:} \textit{“Thưa thầy, Student's t-test đòi hỏi phương sai giữa hai nhóm phải bằng nhau. Trong thực tế đô thị, nồng độ bụi vào giờ cao điểm kẹt xe có độ biến thiên và phương sai lớn hơn rất nhiều so với ban đêm. Welch's t-test hiệu chỉnh bậc tự do theo công thức Welch-Satterthwaite, không đòi hỏi phương sai đồng nhất, giúp kiểm soát sai lầm loại I (Type I error) an toàn và chuẩn xác.”}

\textbf{Hỏi:} \textit{“Kết quả kiểm định của em chứng minh được điều gì?”}\\
\textbf{Đáp:} \textit{“Kết quả cho thấy vào giờ cao điểm, nồng độ PM2.5 trung bình tăng 22.8 ug/m3 so với giờ bình thường, với p-value < 0.001 và t-statistic vượt 35. Điều này bác bỏ hoàn toàn giả thuyết vô hiệu H0, chứng minh kẹt xe làm gia tăng ô nhiễm không khí có ý nghĩa thống kê thực nghiệm.”}
\end{tcolorbox}

% ---------------- TV 5 ----------------
\subsubsection{Phân Hệ 5: Spatial Graph \& Link Analysis (TV 5)}
\textbf{Người phụ trách:} Thành viên 5 \quad $\vert$ \quad \textbf{Tiến độ:} Sprint 3.
\begin{itemize}[leftmargin=*]
    \item \textbf{Đầu vào \& Đầu ra:}
    \begin{itemize}
        \item \textit{Đầu vào:} Cấu trúc tiếp giáp không gian và lưu lượng giao thông liên quận từ TV 3.
        \item \textit{Đầu ra:} Đồ thị NetworkX, bảng xếp hạng điểm nóng lan truyền ô nhiễm theo điểm số PageRank có trọng số.
    \end{itemize}
    \item \textbf{Cơ sở lý thuyết \& Toán học:}
    Thuật toán PageRank có trọng số (Weighted PageRank) trên đồ thị không gian có hướng $G=(V, E)$:
    \begin{equation}
        \text{PR}(u) = \frac{1-d}{|V|} + d \sum_{v \in \text{In}(u)} \frac{W(v, u)}{\sum_{k \in \text{Out}(v)} W(v, k)} \text{PR}(v), \quad (d = 0.85).
    \end{equation}
    Trong đó $W(v, u)$ là trọng số lưu lượng phương tiện và nồng độ ô nhiễm phát thải giữa hai quận tiếp giáp.
    \item \textbf{Lập trình \& Kỹ thuật:} Module \texttt{src/graph/spatial\_graph.py}. Xây dựng đồ thị 5 đỉnh, 12 cạnh có hướng bằng \texttt{networkx}. Tính toán PageRank và Betweenness Centrality.
    \item \textbf{Trực quan hóa \& Báo cáo:} Đồ thị mạng lưới không gian với kích thước Node tỷ lệ thuận với điểm PageRank; Soạn thảo Chương 5 báo cáo.
\end{itemize}

\begin{tcolorbox}[colback=blue!5!white, colframe=blue!60!black, fonttitle=\bfseries\color{white}, title={BỘ CÂU HỎI VẤN ĐÁP TRỌNG TÂM DÀNH CHO TV 5}]
\textbf{Hỏi:} \textit{“Ý tưởng ánh xạ thuật toán PageRank từ xếp hạng web sang xếp hạng ô nhiễm đô thị là gì?”}\\
\textbf{Đáp:} \textit{“Thưa thầy, nếu trong trang web, PageRank đo lường mức độ quan trọng dựa trên các liên kết trỏ đến, thì trong mạng lưới đô thị, một quận đóng vai trò Hub ô nhiễm nếu nó nhận luồng phương tiện và gió mang bụi từ nhiều quận khác đổ về. Thuật toán phân bổ ngẫu nhiên có trọng số giúp xác định chính xác quận nào là điểm nóng tích tụ lan truyền rủi ro cao nhất.”}

\textbf{Hỏi:} \textit{“Kết quả thuật toán chỉ ra quận nào có điểm PageRank cao nhất tại Hà Nội?”}\\
\textbf{Đáp:} \textit{“Kết quả thực nghiệm xác định Cầu Giấy (HN\_CG) và Đống Đa (HN\_DD) đạt điểm PageRank cao nhất (lần lượt là 0.28 và 0.25), vì đây là các nút giao huyết mạch tập trung mật độ phương tiện và kết nối hướng tâm lớn nhất Thủ đô.”}
\end{tcolorbox}

% ---------------- TV 6 ----------------
\subsubsection{Phân Hệ 6: Tabular Machine Learning Predictor (TV 6)}
\textbf{Người phụ trách:} Thành viên 6 \quad $\vert$ \quad \textbf{Tiến độ:} Sprint 3.
\begin{itemize}[leftmargin=*]
    \item \textbf{Đầu vào \& Đầu ra:}
    \begin{itemize}
        \item \textit{Đầu vào:} \texttt{master\_urban\_health.csv} từ TV 3.
        \item \textit{Đầu ra:} Mô hình Hồi quy dự báo nồng độ bụi ($R^2 > 0.40$); Mô hình Phân loại Báo động Đỏ Sức khỏe ($\text{AUROC} \ge 0.75$); Vector xác suất dự báo bàn giao cho TV 8.
    \end{itemize}
    \item \textbf{Cơ sở lý thuyết \& Toán học:}
    Thuật toán Gradient Boosting Regressor tối thiểu hóa hàm mất mát bình phương:
    \begin{equation}
        F_m(x) = F_{m-1}(x) + \gamma_m h_m(x), \quad r_{im} = -\left[\frac{\partial L(y_i, F(x_i))}{\partial F(x_i)}\right]_{F=F_{m-1}}.
    \end{equation}
    Mô hình Phân loại Logistic/Gradient Boosting dự báo xác suất sự cố cực đoan $P(\text{UVSI} \ge 60 \mid X)$. Trích xuất tầm quan trọng đặc trưng (Feature Importance).
    \item \textbf{Lập trình \& Kỹ thuật:} Module \texttt{src/ml\_tabular/health\_predictor.py}. Chia tập dữ liệu theo thứ tự thời gian (80\% Train, 20\% Test). Huấn luyện \texttt{GradientBoostingRegressor} và \texttt{GradientBoostingClassifier} của \texttt{scikit-learn}.
    \item \textbf{Trực quan hóa \& Báo cáo:} Biểu đồ Feature Importance; Biểu đồ Scatter Plot giữa giá trị thực tế và giá trị dự báo; Soạn thảo Chương 6 báo cáo.
\end{itemize}

\begin{tcolorbox}[colback=blue!5!white, colframe=blue!60!black, fonttitle=\bfseries\color{white}, title={BỘ CÂU HỎI VẤN ĐÁP TRỌNG TÂM DÀNH CHO TV 6}]
\textbf{Hỏi:} \textit{“Mô hình Machine Learning của em có ứng dụng thực tế gì hay chỉ dự báo con số trừu tượng?”}\\
\textbf{Đáp:} \textit{“Thưa thầy, mô hình của em giải quyết 3 bài toán thực tế sát sườn: Thứ nhất, tư vấn cho người dân khung giờ tập thể dục an toàn và khuyến cáo đeo khẩu trang N95 trước 3-6 tiếng; Thứ hai, cảnh báo sớm nguy cơ quá tải bệnh nhi cho các bệnh viện hô hấp; Thứ ba, hỗ trợ CSGT mô phỏng chính sách What-If: dự đoán nếu giảm 30\% kẹt xe tại Cầu Giấy thì nồng độ bụi mịn PM2.5 sẽ giảm được bao nhiêu ug/m3.”}

\textbf{Hỏi:} \textit{“Tại sao em dùng phân chia Train/Test theo chuỗi thời gian thay vì K-Fold Cross Validation ngẫu nhiên?”}\\
\textbf{Đáp:} \textit{“Vì dữ liệu môi trường có tính tự tương quan thời gian (temporal autocorrelation). Nếu xáo trộn ngẫu nhiên (random shuffle), dữ liệu tương lai sẽ bị rò rỉ vào tập huấn luyện (Data Leakage), làm mô hình bị lạc quan giả tạo. Phân chia theo thời gian bảo đảm tính trung thực khi triển khai thực tế.”}
\end{tcolorbox}

% ---------------- TV 7 ----------------
\subsubsection{Phân Hệ 7: Incident NLP \& Text Mining (TV 7)}
\textbf{Người phụ trách:} Thành viên 7 \quad $\vert$ \quad \textbf{Tiến độ:} Sprint 3.
\begin{itemize}[leftmargin=*]
    \item \textbf{Đầu vào \& Đầu ra:}
    \begin{itemize}
        \item \textit{Đầu vào:} Tệp tin sự cố văn bản \texttt{vov\_live\_incidents.csv} từ TV 1.
        \item \textit{Đầu ra:} Bảng dữ liệu sự cố đã được bóc tách thực thể địa danh (District Code) và gán nhãn mức độ nghiêm trọng (Severity Impact Score).
    \end{itemize}
    \item \textbf{Cơ sở lý thuyết \& Toán học:}
    Khai phá văn bản (Text Mining), biểu diễn từ ngữ và trích xuất thực thể có tên (Named Entity Recognition -- NER) dựa trên từ điển ngữ cảnh đô thị Hà Nội. Hàm tính điểm tác động rủi ro sự cố:
    \begin{equation}
        \text{Impact\_Score} = w_{\text{type}} \times \mathbb{I}(\text{Incident}) + \beta \times \text{Context\_Weight}.
    \end{equation}
    \item \textbf{Lập trình \& Kỹ thuật:} Module \texttt{src/ml\_nlp/event\_parser.py}. Sử dụng biểu thức chính quy Regex nâng cao kết hợp phân tích từ khóa tiếng Việt (ngập úng, va chạm liên hoàn, mất đèn tín hiệu) để ánh xạ về mã quận tương ứng.
    \item \textbf{Trực quan hóa \& Báo cáo:} WordCloud từ khóa sự cố giao thông; Biểu đồ cột phân loại tai nạn/ngập lụt theo quận; Soạn thảo Chương 7 báo cáo.
\end{itemize}

\begin{tcolorbox}[colback=blue!5!white, colframe=blue!60!black, fonttitle=\bfseries\color{white}, title={BỘ CÂU HỎI VẤN ĐÁP TRỌNG TÂM DÀNH CHO TV 7}]
\textbf{Hỏi:} \textit{“Khi bản tin VOV chỉ nhắc đến tên đường phố (ví dụ: Chùa Bộc, Cầu Giấy) mà không có tên quận, thuật toán của em làm thế nào để định vị?”}\\
\textbf{Đáp:} \textit{“Thưa thầy, em xây dựng một Gazetteer (Từ điển địa danh phân cấp) ánh xạ các trục đường, ngã tư huyết mạch của Hà Nội về quận trực thuộc (ví dụ: Chùa Bộc, Thái Hà -> Đống Đa; Xuân Thủy, Duy Tân -> Cầu Giấy). Nhờ đó, thuật toán nhận diện chính xác quận xảy ra sự cố với độ chính xác cao.”}

\textbf{Hỏi:} \textit{“Phân hệ NLP này đóng góp giá trị gì vào hệ thống chung?”}\\
\textbf{Đáp:} \textit{“Nó cung cấp nguồn dữ liệu động thời gian thực về các biến cố đột xuất (như mưa ngập làm tê liệt giao thông), giúp hệ thống bổ sung tức thì các điểm cảnh báo đột biến vào bản đồ mà các trạm đo định kỳ không giải thích được.”}
\end{tcolorbox}

% ---------------- TV 8 ----------------
\subsubsection{Phân Hệ 8: Model Evaluation \& DeLong ROC Test (TV 8)}
\textbf{Người phụ trách:} Thành viên 8 \quad $\vert$ \quad \textbf{Tiến độ:} Sprint 4.
\begin{itemize}[leftmargin=*]
    \item \textbf{Đầu vào \& Đầu ra:}
    \begin{itemize}
        \item \textit{Đầu vào:} Xác suất dự báo và nhãn thực tế từ mô hình của TV 6.
        \item \textit{Đầu ra:} Báo cáo kiểm định phi tham số DeLong Test, đường cong ROC/PR, giá trị Brier score và ngưỡng phân loại tối ưu Youden's J.
    \end{itemize}
    \item \textbf{Cơ sở lý thuyết \& Toán học:}
    Kiểm định phi tham số DeLong Test so sánh hai diện tích dưới đường cong ROC (AUROC) dựa trên lý thuyết thống kê U-statistics:
    \begin{equation}
        Z = \frac{\widehat{\text{AUC}}_1 - \widehat{\text{AUC}}_2}{\sqrt{\mathbb{V}\text{ar}(\widehat{\text{AUC}}_1 - \widehat{\text{AUC}}_2)}} \sim \mathcal{N}(0, 1).
    \end{equation}
    Tối ưu hóa ngưỡng phân biệt lâm sàng đô thị bằng chỉ số Youden's J:
    \begin{equation}
        J = \max_{c} \{\text{Sensitivity}(c) + \text{Specificity}(c) - 1\}.
    \end{equation}
    \item \textbf{Lập trình \& Kỹ thuật:} Module \texttt{src/evaluation/urban\_evaluator.py}. Tự hiện thực hóa thuật toán DeLong trên nền tảng \texttt{numpy} và \texttt{scipy}. Quét qua 100 ngưỡng xác suất để xác định ngưỡng cắt tối ưu $c^*$.
    \item \textbf{Trực quan hóa \& Báo cáo:} Đồ thị ROC Curve so sánh mô hình Đa thể thức vs Baseline có dải tin cậy 95\%; Biểu đồ Reliability Diagram (Calibration); Soạn thảo Chương 8 báo cáo.
\end{itemize}

\begin{tcolorbox}[colback=blue!5!white, colframe=blue!60!black, fonttitle=\bfseries\color{white}, title={BỘ CÂU HỎI VẤN ĐÁP TRỌNG TÂM DÀNH CHO TV 8}]
\textbf{Hỏi:} \textit{“Bản chất toán học của kiểm định DeLong là gì và tại sao nhóm em sử dụng nó?”}\\
\textbf{Đáp:} \textit{“Thưa thầy, khi mô hình tích hợp có AUC = 0.82 cao hơn mô hình cơ sở AUC = 0.74, sự chênh lệch có thể do ngẫu nhiên lấy mẫu. Kiểm định DeLong sử dụng lý thuyết U-statistics để tính ma trận hiệp phương sai tiệm cận của hai đường cong trên cùng một tập test, cho ra Z-score và p-value chính xác mà không cần giả định phân phối chuẩn, chứng minh sự vượt trội là có ý nghĩa thống kê thực sự.”}

\textbf{Hỏi:} \textit{“Chỉ số Youden's J giúp ích gì cho việc đưa ra quyết định cảnh báo?”}\\
\textbf{Đáp:} \textit{“Ngưỡng mặc định 0.5 trong học máy thường không tối ưu khi dữ liệu mất cân bằng. Youden's J tìm ra điểm cutoff tối ưu cân bằng hoàn hảo giữa độ nhạy (không bỏ sót các đợt bão bụi nguy hại) và độ đặc hiệu (không gây báo động giả làm người dân hoang mang).”}
\end{tcolorbox}

% ---------------- TV 9 ----------------
\subsubsection{Phân Hệ 9: CDSS Web Application \& Master Storytelling (TV 9)}
\textbf{Người phụ trách:} Thành viên 9 \quad $\vert$ \quad \textbf{Tiến độ:} Sprint 4 \& 5 (Tích hợp toàn diện).
\begin{itemize}[leftmargin=*]
    \item \textbf{Đầu vào \& Đầu ra:}
    \begin{itemize}
        \item \textit{Đầu vào:} Sản phẩm, mô hình và dữ liệu của toàn bộ 8 thành viên phía trên.
        \item \textit{Đầu ra:} Cỗ máy hỗ trợ quyết định lâm sàng đô thị (CDSS Engine); Mô phỏng chính sách What-If; Tổng biên tập toàn văn Báo cáo \LaTeX\ và Bộ Master Slide.
    \end{itemize}
    \item \textbf{Cơ sở lý thuyết \& Toán học:}
    Lý thuyết Hệ thống Hỗ trợ Quyết định (Clinical Decision Support Systems -- CDSS). Công thức tính Điểm Đáng Sống Đô Thị (Urban Liveability Score: $0 \to 100$):
    \begin{equation}
        \text{Liveability} = \max\left(0.0, \, 100.0 - 0.5 \times \text{PM2.5} - 4.0 \times \text{Jam\_Factor}\right).
    \end{equation}
    Mô phỏng hàm phản ứng can thiệp chính sách What-If Policy Simulation.
    \item \textbf{Lập trình \& Kỹ thuật:} Module \texttt{src/cdss\_app/urban\_health\_app.py}. Tích hợp giao diện hỗ trợ ra quyết định, thanh trượt can thiệp lưu lượng xe và dự báo mức tăng điểm đáng sống tương ứng.
    \item \textbf{Trực quan hóa \& Báo cáo:} Giao diện điều khiển tương tác CDSS; Tổng duyệt dàn trang báo cáo \LaTeX\ hoàn chỉnh; Soạn thảo Chương 9 báo cáo và Kết luận.
\end{itemize}

\begin{tcolorbox}[colback=blue!5!white, colframe=blue!60!black, fonttitle=\bfseries\color{white}, title={BỘ CÂU HỎI VẤN ĐÁP TRỌNG TÂM DÀNH CHO TV 9}]
\textbf{Hỏi:} \textit{“Mô-đun What-If Policy Simulation của em hoạt động như thế nào?”}\\
\textbf{Đáp:} \textit{“Thưa thầy, khi người dùng hoặc nhà hoạch định chính sách kéo thanh trượt giảm kẹt xe (ví dụ giảm 30\% tại Cầu Giấy), hệ thống sẽ nạp tham số mới vào mô hình dự báo của TV 6 để tính nồng độ PM2.5 kỳ vọng mới, từ đó tính ra số điểm Liveability được cải thiện (tăng từ 52 lên 68 điểm), giúp cơ quan quản lý thấy ngay lợi ích định lượng của quyết sách.”}

\textbf{Hỏi:} \textit{“Em làm thế nào để đảm bảo tính kết nối hữu cơ xuyên suốt giữa 9 thành viên?”}\\
\textbf{Đáp:} \textit{“Em thiết lập luồng dữ liệu khép kín: Dữ liệu thô từ TV 1 -> làm sạch qua TV 2 -> tích hợp Master Table tại TV 3 -> cấp cho TV 4 (thống kê), TV 5 (đồ thị), TV 6 (học máy), TV 7 (NLP) -> thẩm định qua TV 8 -> và hội tụ đầy đủ lên ứng dụng CDSS do em phụ trách.”}
\end{tcolorbox}

% -------------------------------------------------------------
% PHẦN 5: SƠ ĐỒ CHUYỂN GIAO VÀ QUAN HỆ TIÊN QUYẾT
% -------------------------------------------------------------
\section{Sơ Đồ Chuyển Giao \& Quan Hệ Tiên Quyết Giữa 09 Thành Viên}

Quy trình phối hợp và chuyển giao sản phẩm trung gian giữa 9 thành viên được chuẩn hóa theo luồng dây chuyền sản xuất khép kín không chồng chéo:

\begin{center}
\begin{tikzpicture}[auto, >=stealth, font=\footnotesize]
    % Hàng 1: Khai phá & Quản trị Dữ liệu
    \node (tv1) [draw=red!70, fill=red!10, rounded corners, text width=3.8cm, align=center, minimum height=1.4cm] at (0, 0) {\textbf{TV 1 (Data Ingestion)}\\Thu thập Đa nguồn\\Bàn giao: \texttt{data/01\_raw/}};
    \node (tv2) [draw=blue!70, fill=blue!10, rounded corners, text width=3.8cm, align=center, minimum height=1.4cm] at (5.5, 0) {\textbf{TV 2 (Governance)}\\Làm sạch \& Khử nhiễu\\Bàn giao: \texttt{data/02\_curated/}};
    \node (tv3) [draw=forestgreen!70, fill=forestgreen!10, rounded corners, text width=3.8cm, align=center, minimum height=1.4cm] at (11, 0) {\textbf{TV 3 (Integration)}\\Spatiotemporal Join\\Bàn giao: \texttt{Master Dataset}};
    
    % Hàng 2: Mô hình hóa & Phân tích từ Master Table
    \node (tv4) [draw=orange!70, fill=orange!10, rounded corners, text width=3.5cm, align=center, minimum height=1.3cm] at (0, -2.8) {\textbf{TV 4 (EDA / Thống kê)}\\Kiểm định Welch t-test\\$p < 0.001$};
    \node (tv5) [draw=teal!70, fill=teal!10, rounded corners, text width=3.5cm, align=center, minimum height=1.3cm] at (5.5, -2.8) {\textbf{TV 5 (Graph Analytics)}\\Spatial PageRank\\Hub ô nhiễm};
    \node (tv6) [draw=purple!70, fill=purple!10, rounded corners, text width=3.5cm, align=center, minimum height=1.3cm] at (11, -2.8) {\textbf{TV 6 (Tabular ML)}\\Gradient Boosting\\Dự báo $\text{PM}_{2.5}$};
    
    % Hàng 3: Xử lý văn bản, Thẩm định & Dashboard CDSS
    \node (tv7) [draw=brown!70, fill=brown!10, rounded corners, text width=3.5cm, align=center, minimum height=1.3cm] at (0, -5.6) {\textbf{TV 7 (Incident NLP)}\\Bóc tách sự cố\\NER VOV News};
    \node (tv8) [draw=magenta!70, fill=magenta!10, rounded corners, text width=3.5cm, align=center, minimum height=1.3cm] at (5.5, -5.6) {\textbf{TV 8 (Evaluation)}\\Kiểm định DeLong\\Youden's J Cutoff};
    \node (tv9) [draw=navyblue!70, fill=navyblue!10, rounded corners, text width=3.5cm, align=center, minimum height=1.3cm] at (11, -5.6) {\textbf{TV 9 (CDSS Web App)}\\Điểm Đáng sống \&\\Mô phỏng What-If};

    % Đường kết nối Hàng 1
    \draw[->, line width=1.1pt] (tv1) -- (tv2);
    \draw[->, line width=1.1pt] (tv2) -- (tv3);
    
    % TV 3 cấp dữ liệu xuống TV 4, TV 5, TV 6
    \draw[->, line width=1.1pt] (tv3.south) -- (tv6.north);
    \draw[->, line width=1.1pt] (tv3.south) -- ++(0,-0.6) -| (tv5.north);
    \draw[->, line width=1.1pt] (tv3.south) -- ++(0,-0.6) -| (tv4.north);
    
    % TV 1 cấp tin tức VOV xuống TV 7 (chạy vòng bên mép trái)
    \draw[->, line width=1.1pt] (tv1.west) -- ++(-0.6,0) |- (tv7.west);
    
    % TV 6 cấp xác suất dự báo sang TV 8
    \draw[->, line width=1.1pt] (tv6.south) -- ++(0,-0.6) -| (tv8.north);
    
    % TV 6 cấp mô hình sang TV 9
    \draw[->, line width=1.1pt] (tv6.south) -- (tv9.north);
    
    % TV 8 cấp kết quả thẩm định sang TV 9
    \draw[->, line width=1.1pt] (tv8.east) -- (tv9.west);
    
    % TV 7 cấp sự cố sang TV 9 (chạy vòng mép dưới)
    \draw[->, line width=1.1pt] (tv7.south) -- ++(0,-0.6) -| (tv9.south);
    
    % TV 4 & TV 5 cấp kết quả sang TV 9
    \draw[->, line width=1.1pt] (tv5.south) -- ++(0,-0.4) -| (tv9.north);
\end{tikzpicture}
\end{center}

% -------------------------------------------------------------
% PHẦN 6: LỘ TRÌNH TRIỂN KHAI 15 TUẦN
% -------------------------------------------------------------
\section{Lộ Trình Triển Khai Chi Tiết (06 Sprints / 15 Tuần)}

Tiến độ dự án được quản trị nghiêm ngặt theo mô hình Agile/Scrum với 06 Sprints trải dài trong 15 tuần học kỳ:

\begin{table}[htbp]
\centering
\small
\renewcommand{\arraystretch}{1.3}
\caption{Kế hoạch tiến độ triển khai 15 tuần của dự án IT4142.}
\label{tab:sprint_plan}
\begin{tabularx}{\textwidth}{>{\centering\arraybackslash}p{1.6cm} >{\centering\arraybackslash}p{2.0cm} >{\raggedright\arraybackslash}X >{\raggedright\arraybackslash}p{4.6cm}}
\toprule
\textbf{Sprint} & \textbf{Thời gian} & \textbf{Mục tiêu trọng tâm \& Nội dung công việc} & \textbf{Sản phẩm bàn giao} \\
\midrule
\textbf{Sprint 1} & Tuần 1 -- 3 & Thống nhất đề tài; TV 1 xây dựng crawler đa nguồn; TV 2 thiết lập pipeline làm sạch, lọc mã lỗi -999, xuất Data Docs HTML. & \texttt{data/01\_raw/},\newline \texttt{data/02\_curated/},\newline Data Docs HTML. \\
\textbf{Sprint 2} & Tuần 4 -- 6 & TV 3 hoàn thiện ghép bảng Spatiotemporal Join đạt $>10.000$ dòng; TV 4 tiến hành kiểm định Welch's t-test và phân tích tương quan. & \texttt{master\_urban\_health.csv},\newline Báo cáo thống kê Table 1. \\
\textbf{Sprint 3} & Tuần 7 -- 9 & TV 5 xây dựng đồ thị PageRank; TV 6 huấn luyện mô hình Gradient Boosting; TV 7 bóc tách tin tức VOV qua NER. & Điểm PageRank Hotspot,\newline Trọng số mô hình ML,\newline File bóc tách thực thể. \\
\textbf{Sprint 4} & Tuần 10 -- 11 & TV 8 thực hiện kiểm định DeLong ROC và tối ưu ngưỡng Youden's J; TV 9 phát triển khung giao diện CDSS và điểm Liveability. & Kết quả kiểm định DeLong,\newline Bản Alpha CDSS App. \\
\textbf{Sprint 5} & Tuần 12 & Mỗi thành viên hoàn thành viết chương báo cáo cá nhân bằng \LaTeX; thiết kế 3--4 slide cá nhân theo Master Template; Tổng biên tập. & Bản dự thảo toàn văn Báo cáo \LaTeX\ và Bộ Slide thuyết trình. \\
\textbf{Sprint 6} & Tuần 13 -- 15 & Diễn tập vấn đáp giả lập (Mock Defense) 02 vòng nội bộ; hoàn thiện mã nguồn, đóng gói tài liệu và bảo vệ chính thức trước Thầy. & Repository chuẩn mực, Toàn đội bảo vệ đạt điểm tuyệt đối A/A+. \\
\bottomrule
\end{tabularx}
\end{table}

% -------------------------------------------------------------
% PHẦN 7: QUY CHUẨN ĐỒ HỌA & TRỰC QUAN HÓA
% -------------------------------------------------------------
\section{Quy Chuẩn Đồ Họa \& Trực Quan Hóa Đô Thị Chuẩn Mực}

Nhằm đáp ứng yêu cầu khắt khe về tính học thuật và thẩm mỹ trực quan (dựa trên nguyên tắc của Claus Wilke và Edward Tufte), nhóm tuân thủ triệt để các tiêu chí sau:
\begin{enumerate}
    \item \textbf{Quy chuẩn Bảng Màu Đô Thị (Urban Environmental Palette):}
    \begin{itemize}
        \item Nghiêm cấm sử dụng các màu sắc gây cảm xúc tiêu cực hoặc báo động giả. Nhóm điều kiện không khí an toàn sử dụng màu xanh lam sẫm (\texttt{\#2B4C7E}); nhóm ô nhiễm nguy hại sử dụng màu đỏ gạch cảnh báo (\texttt{\#D9534F}).
        \item Ma trận tương quan bắt buộc sử dụng Diverging Palette (\texttt{coolwarm} hoặc \texttt{vlag}), với mức 0 nằm chính xác tại điểm trung hòa màu trắng.
        \item Tuyệt đối không sử dụng bảng màu cầu vồng (\textit{Rainbow / Jet palette}) vì hiện tượng méo mó nhận thức gradient.
    \end{itemize}
    \item \textbf{Tối ưu Hóa Tỷ Lệ Data-to-Ink (Edward Tufte):}
    \begin{itemize}
        \item Loại bỏ các đường kẻ viền hộp không cần thiết, loại bỏ lưới nền đậm màu; chỉ giữ lại đường gióng nhạt ở các ngưỡng tiêu chuẩn EPA ($\text{PM2.5} = 35.4$ và $55.4\,\mu\text{g/m}^3$).
        \item Gán nhãn trực tiếp (Direct Labeling) tại điểm cuối của đường cong thay vì dùng hộp chú giải (legend) ở góc xa.
    \end{itemize}
    \item \textbf{Trực quan hóa Sự Không Chắc Chắn (Visualizing Uncertainty):}
    \begin{itemize}
        \item Không bao giờ trình bày một con số ước lượng điểm đơn độc. Toàn bộ các đồ thị nồng độ bụi theo giờ, đường cong phản ứng phi tuyến và đường cong ROC bắt buộc phải đi kèm dải bóng mờ biểu diễn khoảng tin cậy 95\% (\textit{95\% Shaded Confidence Intervals}).
    \end{itemize}
\end{enumerate}

% -------------------------------------------------------------
% PHẦN 8: CHIẾN THUẬT VẤN ĐÁP 3 BƯỚC
% -------------------------------------------------------------
\section{Chiến Thuật Ghi Điểm Tuyệt Đối Trong Phiên Vấn Đáp}

Mỗi thành viên trong nhóm rèn luyện kỹ năng phản biện theo Chiến thuật vấn đáp 3 bước chuẩn mực:

\begin{tcolorbox}[colback=red!5!white, colframe=red!60!black, fonttitle=\bfseries\color{white}, title={CHIẾN THUẬT VẤN ĐÁP 3 BƯỚC CHO MỖI THÀNH VIÊN (120 GIÂY)}]
\begin{enumerate}[leftmargin=*]
    \item \textbf{Bước 1: Nêu bật bản chất toán học / thống kê (30 giây):}\\
    \textit{“Thưa thầy, cơ sở lý thuyết cốt lõi trong phân hệ của em là [Công thức / Định lý / Thuật toán], nhằm giải quyết trực tiếp câu hỏi nghiệp vụ đô thị là [...]”}
    \item \textbf{Bước 2: Giải trình dòng lệnh mã nguồn \& Tối ưu hóa kỹ thuật (45 giây):}\\
    \textit{“Em đã trực tiếp lập trình giải thuật này tại tệp tin [tên\_tệp.py], áp dụng [thư viện chuyên sâu / thuật toán tối ưu] và giải quyết bài toán hiệu năng ở điểm [...]”}
    \item \textbf{Bước 3: Diễn giải ý nghĩa biểu đồ \& Giá trị thực tiễn (45 giây):}\\
    \textit{“Kết quả trực quan thu được thể hiện trên biểu đồ [Tên biểu đồ]. Con số này không đơn thuần là tham số kỹ thuật, mà chứng minh cho người dân và cơ quan quản lý đô thị thấy rằng [...]”}
\end{enumerate}
\end{tcolorbox}

\begin{itemize}
    \item \textbf{Trình diễn Slide mở đầu:} Slide đầu tiên của nhóm sẽ thể hiện minh bạch bảng phân công và tỷ lệ đóng góp $11.11\% \times 9$, khẳng định tính độc lập khép kín của từng thành viên.
    \item \textbf{Bản lĩnh học thuật:} Khi giảng viên hỏi mở rộng liên quan đến phân hệ của thành viên khác, người được hỏi vẫn có thể trình bày mạch lạc bức tranh tổng thể và điểm giao tiếp (interface contract) giữa phần của mình và phần của đồng đội.
\end{itemize}

% -------------------------------------------------------------
% PHẦN 9: KẾT LUẬN & TÀI LIỆU THAM KHẢO
% -------------------------------------------------------------
\section{Kết Luận}
Bản kế hoạch triển khai Bài tập lớn IT4142 trên đây đã chuẩn hóa đầy đủ cơ sở lý thuyết, kiến trúc công nghệ Lakehouse, phân rã trách nhiệm và chiến lược bảo vệ cho nhóm 09 sinh viên. Với sự chuẩn bị nghiêm túc, bám sát các công nghệ quản trị dữ liệu mới nhất (Polars, Parquet, Data Docs HTML, NetworkX, Scikit-learn, DeLong Test) cùng quy trình phân tích chuỗi thời gian thực nghiệm chặt chẽ, nhóm tự tin hướng tới mục tiêu đạt kết quả cao nhất (Điểm A/A+) trong học phần Nhập môn Khoa học Dữ liệu.

\begin{thebibliography}{99}
\bibitem{hastie2009} T. Hastie, R. Tibshirani, and J. Friedman, \textit{The Elements of Statistical Learning: Data Mining, Inference, and Prediction}, 2nd ed., Springer, New York, NY, 2009.
\bibitem{delong1988} E. R. DeLong, D. M. DeLong, and D. L. Clarke-Pearson, “Comparing the areas under two or more correlated receiver operating characteristic curves: A nonparametric approach,” \textit{Biometrics}, vol. 44, no. 3, pp. 837--845, 1988.
\bibitem{wilke2019} C. O. Wilke, \textit{Fundamentals of Data Visualization: A Primer on Making Informative and Compelling Figures}, O'Reilly Media, Sebastopol, CA, 2019.
\bibitem{tufte2001} E. R. Tufte, \textit{The Visual Display of Quantitative Information}, 2nd ed., Graphics Press, Cheshire, CT, 2001.
\bibitem{page1999} L. Page, S. Brin, R. Motwani, and T. Winograd, “The PageRank citation ranking: Bringing order to the web,” Stanford InfoLab Technical Report, 1999.
\bibitem{knaflic2015} C. N. Knaflic, \textit{Storytelling with Data: A Data Visualization Guide for Business Professionals}, John Wiley \& Sons, Hoboken, NJ, 2015.
\bibitem{armbrust2020} M. Armbrust et al., “Delta Lake: High-performance ACID table storage over cloud object stores,” in \textit{Proceedings of the VLDB Endowment}, vol. 13, no. 12, pp. 3411--3424, 2020.
\bibitem{ritchie2023} Ritchie Vink et al., “Polars: Blazingly fast DataFrames in Rust and Python,” \url{https://pola.rs}, 2023.
\end{thebibliography}

\end{document}
